\section{实验记录的保存位置}
论文源码目录仅保留 \nolinkurl{bib/}、\nolinkurl{pic/}、\nolinkurl{chapter/} 和
\nolinkurl{main.tex}。表格与统计值已经写入相应章节，编译不需要 Python、
检测器、原始样本库或实验数据目录。

实验记录仍有核验和复现用途，已移出论文源码目录。
以下路径均相对于工作区根目录；整理前的完整目录位于
\nolinkurl{essay_support/2026-10-16/original/}，下文简称 R。
主工具及产品示例仍分别位于 \nolinkurl{ustc_tool_middle/} 和 \nolinkurl{examples/}。
原始报告中的输入路径是生成报告时的来源标识，整理不改变其中的内容哈希或实验计数。

\begin{description}
  \item[E01] 完整开发结果：
  \nolinkurl{ustc_tool_middle/build/reports/miasm-cfg-20261015/development-final-summary.json}。
  保存 2344 份输入的真值、已知范围、预测、扫描状态和分算法结果。

  \item[E02] 历史验证复查：同目录下
  \nolinkurl{history-summary.json}、\nolinkurl{recent-summary.json}、\nolinkurl{rf-summary.json}，
  分别对应 2855、232、252 份输入。

  \item[E03] Miasm 对照：同目录下 \nolinkurl{validation-comparison.json}，
  记录接入前后预测、原语、候选状态与阶段耗时。

  \item[E04] 开发对照与冻结：同目录下
  \nolinkurl{development-final-comparison.json} 和 \nolinkurl{final-verification.json}，
  记录 560 份受影响重扫、1784 份复用以及输入完整性。

  \item[E05] 依赖与回归：同目录下
  \nolinkurl{doctor-final.json}、\nolinkurl{regression-final-summary.json}，
  保留依赖、签名库版本及 240 方法、28 模块的真实回归状态。

  \item[E06] 原二进制演示：
  \nolinkurl{ustc_tool_middle/build/reports/competition-selection-20261015/}，
  保留 437 份筛选输入的复扫、架构分组和选择说明。

  \item[E07] 虚拟网关：\nolinkurl{examples/gridlens/verification.json} 和
  \nolinkurl{examples/gridlens/data/demo.json}，
  记录四视图资产、11 项检查、浏览器与现场任务结果。

  \item[E08] 图片来源：R 下 \nolinkurl{pic/sources.json}、\nolinkurl{pic/README.md}，
  保留图片内容哈希、裁切、原始数据和来源说明；当前源码的
  \nolinkurl{pic/} 只保存正文实际使用的 19 张 PNG。
  P18/P19 为第一章融合时依据文献补绘的时间示意，绘图脚本与来源说明位于
  \nolinkurl{essay_support/2026-10-16/chapter1-merge/}。

  \item[E09] 五算法混合演示：R 下 \nolinkurl{showcase/20261016/}，
  保存清单、内容锁、输入与上下文、241 份实际扫描报告及来源分组。
  其紧凑结果副本在 R 下 \nolinkurl{data/showcase-20261016/}。

  \item[E10] 五类终端固件：\nolinkurl{examples/terminal-lab/} 下的
  \nolinkurl{manifests/}、\nolinkurl{firmware/}、\nolinkurl{extracted/}、
  \nolinkurl{data/raw/} 和 \nolinkurl{evidence/}，
  保存 10 份镜像、实际解包、20 个 ELF 报告、五类业务运行与浏览器证据。
  论文使用的结果副本在 R 下 \nolinkurl{data/terminal-lab-20261016/}。

  \item[E11] 十三算法：R 下 \nolinkurl{algorithm-suite/20261016/}，
  保存完整总体、纳入清单、输入与上下文内容锁，以及
  \nolinkurl{fresh-scan/} 中的冻结、1071 份实际新扫报告和结果。
  可脱离检测器重算预测的报告副本在 R 下 \nolinkurl{data/algorithm-suite-20261016/}。
\end{description}

E01、E02、E06 的表格已按每行的已知范围重新核算，计数与原汇总一致。
E09、E10、E11 分别保留了实际扫描报告；三个实验的输入和分母不合并。
当前正文只在显示时将百分比四舍五入到两位小数，
原始整数计数与哈希留在上述记录中。

\section{编译入口与原实验运行器}
本源码使用用户指定的 \nolinkurl{ctexbook} 文档类与宏包，
采用左右 3.5 cm、上 4 cm、下 5 cm 的页边距。
参考文献仅保留实际引用的条目，集中在 \nolinkurl{bib/references.bib}。
建议使用具备中文字体及所列宏包的 XeLaTeX 与 BibTeX，
也可使用 Tectonic 自动处理引用与交叉引用。

\begin{lstlisting}[basicstyle=\ttfamily\footnotesize,breaklines=true,caption={在论文源码目录编译}]
xelatex main.tex
bibtex main
xelatex main.tex
xelatex main.tex
\end{lstlisting}

若要保持源码目录整洁，可以把编译输出指定到外部目录。
本轮实际编译使用 Tectonic，PDF、日志和核验结果保存在
\nolinkurl{essay_support/2026-10-16/build/}，不放入论文源码目录。
当前只有一个 \nolinkurl{main.tex} 入口，不再依赖旧的正文/批注双入口。
配图建议作为章节中的普通 LaTeX 注释保留，不在正文打印。

原实验脚本位于 R 下 \nolinkurl{tools/}。
这些脚本记录整理前的数据生成与扫描流程，其路径约定对应整理前布局，
重新使用时应在原目录结构的工作副本中恢复并配置主项目路径；
不要把它们当作当前纯净论文目录的编译依赖。
主检测器入口仍在 \nolinkurl{ustc_tool_middle/}，
五类终端构建、解包、扫描和核验入口仍在 \nolinkurl{examples/terminal-lab/tools/}。

\section{图像编排与解释边界}
P01--P19 是素材编号，正文图号由 LaTeX 按章生成。
P01/P02 是真实文件内容排版图，P03 是系统架构图；
P04、P08--P10、P14/P15 是实际软件界面，
P05--P07、P11--P13、P16/P17 为从已完成实验记录生成的统计图。
P16 显示十三算法的检出数、漏报数、正标签分母与 10\% 上限；
P17 分别显示源码、x86-64、ARM64 的漏报率和 FN/正标签数。
P18/P19 分别展示先存后解与后量子迁移节点的时间关系，属于文献示意，不是实测结果。
图片没有将厂商照片、实体设备或未执行实验冒充本项目现场。

本报告仍未完成新的未知家族盲测、完整多通道消融、其他工具同条件对比、
实体电力终端、完整网关运行、后量子协议互操作或原生 Windows 宿主验收。
五类固件是项目自制的最小业务文件系统，ARM64 没有业务运行证据；
十三算法实验是已知开发家族的范围内回归。
历史未知实现覆盖差异与原有回归失败继续按实际状态保留。
